\clearpage\pagestyle{otMain}\fontsize{11}{15}\selectfont
\corrigesChapitre{7}{Variations et optimisation}
\begin{blocLivre}{cor-7-1}
\begin{corrige}{1}{Une pente, un sens, un signe}
\demarche{Repérer la pente pour les variations, puis la valeur qui annule la fonction pour le signe.}
Le signe du coefficient directeur suffit : 4 est positif, donc $f$ est strictement croissante ; $-3$ est négatif, donc $g$ est strictement décroissante ; $h$ a un coefficient nul, donc elle est constante. Le terme ajouté ne change pas le sens de variation. Pour les signes, $f$ s’annule en $\res{7/4}$ : $f(x)$ est négatif pour $x<7/4$ et positif pour $x>7/4$. La fonction $g$ s’annule en $\res{2/3}$ : $g(x)$ est positif pour $x<2/3$ et négatif pour $x>2/3$. Une fonction affine croissante passe du négatif au positif, une décroissante fait l’inverse ; $h$ reste positive.
\end{corrige}
\end{blocLivre}
\begin{blocLivre}{cor-7-2}
\begin{corrige}{2}{Lire une ligne brisée}
\demarche{Parcourir la courbe de gauche à droite.}
La courbe monte de $-2$ à 0, descend de 0 à 2, puis remonte de 2 à 4. Le maximum est $\res3$, atteint en $\res{x=0}$. Le minimum est $\res{-1}$, atteint en $\res{x=2}$. Les valeurs aux deux autres extrémités valent 1 et ne dépassent pas ces extrémums.
\label{fig-cor-7-2}
\figureGeo[corrige]{fonction-morceaux}
\end{corrige}
\end{blocLivre}
\begin{blocLivre}{cor-7-3}
\begin{corrige}{3}{Carré sur un intervalle négatif}
\demarche{Choisir la bonne branche de la fonction carré.}
La fonction carré est décroissante sur les réels négatifs. Sur $[-5;-1]$, elle descend de 25 à 1. Le minimum est $\res1$ en $x=-1$ ; le maximum est $\res{25}$ en $x=-5$. La plus petite abscisse donne ici la plus grande image.
\end{corrige}
\end{blocLivre}
\begin{blocLivre}{cor-7-4}
\begin{corrige}{4}{Un carré décalé}
\demarche{Utiliser la positivité du carré et son cas d’annulation.}
Le carré $(x-2)^2$ est toujours positif ou nul. On en déduit $f(x)\ge3$. L’égalité est possible et exige $x-2=0$. Donc le minimum est $\res3$, atteint uniquement en $\res{x=2}$. Il n’existe pas de maximum sur $\mathbb R$, car le carré peut être arbitrairement grand.
\end{corrige}
\end{blocLivre}
\begin{blocLivre}{cor-7-5}
\begin{corrige}{5}{Comparer sans connaître les images}
\demarche{Séparer comparaison d’images et comparaison à zéro.}
Comme $2<7$ et que la fonction décroît strictement, $\res{f(2)>f(7)}$. On ne peut rien conclure sur le signe de $f(2)$. Les fonctions $-x$ et $10-x$ sont toutes deux décroissantes, mais elles donnent en 2 les valeurs $-2$ et 8, de signes opposés.
\end{corrige}
\end{blocLivre}
\begin{blocLivre}{cor-7-6}
\begin{corrige}{6}{Un tableau ne donne pas tout}
\demarche{Utiliser uniquement les comparaisons garanties par les flèches.}
La croissance donne $\res{1\le f(1)\le5}$. La décroissance sur $[2;6]$ donne $\res{-3\le f(4)\le5}$. Ces valeurs ne sont pas déterminées exactement : la fonction peut suivre de nombreuses courbes entre les extrémités, tout en conservant les mêmes variations. Une interpolation linéaire nécessiterait une hypothèse supplémentaire.
\label{fig-cor-7-6}
\figureGeo[corrige]{variations-lecture}
\end{corrige}
\end{blocLivre}
\begin{blocLivre}{cor-7-7}
\begin{corrige}{7}{Prouver la croissance d’une affine}
\demarche{Démontrer la monotonie par le signe d’une différence d’images.}
On calcule $f(b)-f(a)=(3b/2-4)-(3a/2-4)=\tfrac32(b-a)$. Or $b-a>0$ et $3/2>0$. La différence est donc strictement positive. Ainsi $f(a)<f(b)$ pour tous $a<b$, ce qui est exactement la définition de la croissance stricte.
\end{corrige}
\end{blocLivre}
\begin{blocLivre}{cor-7-8}
\begin{corrige}{8}{Une symétrie utile}
\demarche{Repérer le centre du carré puis comparer les deux extrémités.}
La distance de $x$ à 3 diminue jusqu’à 3 puis augmente. Le carré suit ces variations : $f$ décroît sur $[0;3]$ et croît sur $[3;7]$. Le minimum est $\res{-2}$ en 3. Aux bornes, $f(0)=7$ et $f(7)=14$ ; le maximum sur l’intervalle est donc $\res{14}$ en 7.
\label{fig-cor-7-8}
\figureGeo[corrige]{variation-decalee}
\end{corrige}
\end{blocLivre}
\begin{blocLivre}{cor-7-9}
\begin{corrige}{9}{Maximum d’un produit}
\demarche{Faire apparaître un nombre fixe moins un carré.}
On complète le carré : $A(x)=10x-x^2=25-(x-5)^2$. Le carré est positif ou nul, donc $A(x)\le25$. L’égalité se produit seulement en $x=5$, qui appartient au domaine. Ainsi $\res{\max A=25}$, atteint en $\res5$.
\label{fig-cor-7-9}
\figureGeo[corrige]{aire-produit}
\end{corrige}
\end{blocLivre}
\begin{blocLivre}{cor-7-10}
\begin{corrige}{10}{Un balayage et sa limite}
\demarche{Distinguer optimisation discrète et optimisation sur un intervalle.}
Les images sont 2,31 ; 3,91 ; 3,51 ; 1,11. La grille donne $\res{3{,}91}$ en 1. Mais l’expression montre que le vrai maximum est $\res4$ en $x=1{,}3$. La meilleure valeur d’une grille est seulement un maximum parmi les points testés ; le pas peut manquer l’optimum réel.
\end{corrige}
\end{blocLivre}
\begin{blocLivre}{cor-7-11}
\begin{corrige}{11}{Un enclos contre un mur}
\demarche{Écrire la contrainte de périmètre avant la fonction d’aire.}
Les deux largeurs utilisent $2x$ mètres ; la longueur vaut $24-2x$. Le domaine est $0<x<12$. L’aire vaut $A(x)=x(24-2x)=72-2(x-6)^2$. Elle ne dépasse pas $\res{72\text{ m}^2}$, atteints pour $\res{x=6\text{ m}}$ et une longueur de $\res{12\text{ m}}$. Ces dimensions utilisent bien $6+12+6=24$ m de clôture.
\label{fig-cor-7-11}
\figureGeo[corrige]{enclos}
\end{corrige}
\end{blocLivre}
\begin{blocLivre}{cor-7-12}
\begin{corrige}{12}{La recette d’une vente}
\demarche{Modéliser le produit puis établir un majorant atteint.}
La recette est prix fois quantité : $R(p)=p(100-5p)=500-5(p-10)^2$. Le maximum est donc $\res{500\text{ euros}}$, obtenu pour $\res{p=10\text{ euros}}$. Le modèle prévoit alors 50 ventes. Le résultat dépend de l’hypothèse de demande donnée ; ce calcul ne suffit pas à valider le modèle dans une situation réelle.
\end{corrige}
\end{blocLivre}
\begin{blocLivre}{cor-7-13}
\begin{corrige}{13}{Un minimum de distance au carré}
\demarche{Optimiser le carré de la distance pour éviter la racine.}
La formule de distance donne $AM^2=(x-2)^2+9\ge9$. L’égalité est atteinte seulement pour $x=2$. La racine carrée étant croissante sur les positifs, minimiser $AM^2$ revient à minimiser $AM$. Le point cherché est $\res{M(2;0)}$, et la distance minimale vaut $\res3$.
\label{fig-cor-7-13}
\figureGeo[corrige]{distance-minimale}
\end{corrige}
\end{blocLivre}
\begin{blocLivre}{cor-7-14}
\begin{corrige}{14}{Un maximum sur un domaine restreint}
\demarche{Vérifier l’appartenance du candidat au domaine.}
La valeur 9 serait atteinte en 4, mais 4 n’appartient pas au domaine. Sur $[0;2]$, la distance à 4 diminue : $f$ est croissante. On calcule $f(0)=-7$ et $f(2)=5$. Le minimum est $\res{-7}$ en 0 et le maximum $\res5$ en 2. Un majorant non atteint n’est pas le maximum recherché.
\end{corrige}
\end{blocLivre}
\begin{blocLivre}{cor-7-15}
\begin{corrige}{15}{Une garantie pour toutes les valeurs}
\demarche{Démontrer une borne uniforme puis exhiber une valeur qui l’atteint.}
On écrit $x(8-x)=16-(x-4)^2\le16$. La valeur $k=16$ convient donc pour toutes les abscisses. Si $k<16$, le choix $x=4$ donne $x(8-x)=16>k$, et l’inégalité échoue. Le plus petit majorant demandé est ainsi $\res{k=16}$. On a prouvé à la fois qu’il convient et qu’on ne peut l’abaisser.
\end{corrige}
\end{blocLivre}
